\subsection{The marked symplectic induction in Case II}
\label{sec:rh-symplectic-induction}

\begin{theorem}[Integral rank-four seed in the second case]
\label{thm:rh-seed-II-alt}
Keep the original datum \(D\), the odd coefficient prime \(\ell\),
and the tower \(T\), with precisely the hypotheses
\(\texttt{Ctx2}(D)\), \(\texttt{Ok}(T)\), \(\texttt{Two}(T)\), and
\(\texttt{Deep}(T)\). For every prescribed coefficient embedding
\(\iota:\overline{\mathbb Q}_{\ell}\simeq\mathbb C\), there is an
extension of the coefficient DVR and a pair on the original frame
\(T.\mathrm{fr}(\varnothing)\) satisfying \(\texttt{SeedPairAlt}\),
\(\texttt{InV}\) after reduction, and the original alternative
\[
 \operatorname{Ent}(-w_{\rm Al})\quad\text{or}\quad
 \bigl(\operatorname{Ent}(w_{\rm Be})\ \text{and}\
       \operatorname{Ent}(-w_{\rm Th})\bigr).
\]
In particular, this is the full supplier \(\texttt{SeedIIAlt}\), using
the mixed Hilbert seed of Section~\ref{sec:mixed-hilbert-seed} and
the classical theorems specified below. It is not an assertion of
automorphy for every point on the tower's deformation space.
\end{theorem}

\begin{proof}
We identify all the absolute Galois groups with subgroups of
\(G_{\mathbb Q}\) using the original algebraic closure. Write
\(F=T.F,\ L=T.L=FK\), and \(s=T.g_0\in G_F\), where \(K\) is
the original imaginary quadratic field. Retain every original local
group, denoted \(D_v\) and \(I_v\) above \(\ell\). Write
\(\epsilon\) for the cyclotomic character, \(\omega\) for its
Teichmuller part, \(\xi\) for the finite odd dyadic character of
\(D\), and put
\[
 \alpha_0=\epsilon\omega^{-1}\xi,\qquad
 \Psi=\epsilon^e\psi,\qquad \nu=\alpha_0\Psi.
\]
The ordered frame characters \(\chi_i\), \(0\le i\le3\), are
\((\nu,\Psi,\alpha_0,1)\) or \((\alpha_0,1,\nu,\Psi)\)
according to the original ordering parameter. In both orderings
\(\chi_0\chi_3=\chi_1\chi_2=\nu\), and
\(\bar\alpha_0=\bar\xi\).

\paragraph{Admission of the Hilbert seed.}
By \(\texttt{Two.quad}\) there is the specified real quadratic
field \(k_0\), split at \(2,\ell\). Put \(H=G_{k_0}\) and
let \(\phi:G_{\mathbb Q}\to\{1,-1\}\) be its quadratic character.
On \(G_F\), \(\phi=\bar\Psi\) in the residue field: the ratio
defining \(w_{\rm Ps}\) is \(\Psi\) or \(\Psi^{-1}\), and its
reduction has square one. Its kernel there is \(G_F\cap H\).
It is unramified on all relative inertia groups of \(F\).
By \(\texttt{Ok.dist}\), the four residual characters at \(s\)
are distinct. Hence \(\phi(s)=-1\), \(\bar\xi(s)\ne\pm1\).
In particular \(\bar\xi(s^2)\ne1\). On \(G_F\cap H\) and
on the local groups at \(\ell\), the ordered reductions are
\((\bar\xi,1,\bar\xi,1)\).

The hypotheses \(\texttt{Two.ord}\), \(\texttt{Two.xi}\),
\(\texttt{Two.big2}\), \(\texttt{Two.bigL}\) show, respectively,
that \(\xi\) has \(2\)-power order, is unramified outside \(2\),
has residual order greater than four on a relative inertia group
at \(2\), and has nontrivial residual value on every decomposition
group at \(\ell\). Here the root \(w_{\rm Al}\) has residual
character \(\bar\xi\). Since \(2,\ell\) split in \(k_0\), these
statements apply to \(\xi|_H\) at the places of norm \(2\) and
at each of the two places over \(\ell\). The restriction is totally
odd. Thus the mixed Hilbert theorem applies with one of the two
\(\ell\)-adic places selected. It gives an absolutely irreducible
representation \(r\) of \(H\), over a finite extension of
\(\mathbb Q_\ell\), irreducible on every open subgroup, satisfying
\[
 \det r=\epsilon^k\xi|_H,\qquad
 \bar r^{\rm ss}=1\oplus\bar\xi|_H,
 \qquad k,t\ge2,\quad \ell-1\mid k,t.
\]
At the two \(\ell\)-adic places its ordered inertia exponents are
\((k+t,-t)\) and \((k,0)\), and the ordered residual characters on
the whole decomposition groups are respectively
\((\bar\xi,1)\) and \((1,\bar\xi)\). The conjugation by \(s\)
interchanges these places. Away from \(\ell\) its inertia
representation is exactly \(1\oplus\xi\).

\paragraph{A single marked integral lattice.}
Enlarge the coefficient field to contain the preceding
representation and the original coefficients; call its ring of
integers \(\mathcal O\), maximal ideal \(\mathfrak m\), and
uniformizer \(\varpi\). Compactness gives a stable lattice for
\(H\). The eigenvalues of \(r(s^2)\) are distinct modulo
\(\mathfrak m\), reducing to \(1,\bar\xi(s^2)\). Its integral
spectral idempotents therefore split that lattice into two
rank-one summands, giving an eigenbasis in this order. For
\(h\in H_F=H\cap G_F\), write its matrix in that basis as
\(\left(\begin{smallmatrix}d_h&b_h\\c_h&f_h\end{smallmatrix}\right)\).
The traces at \(h\) and \(s^2h\) give
\(d_h\equiv1,\ f_h\equiv\xi(h)\pmod{\mathfrak m}\).
Multiplication at \(h,h'\) now gives
\(b_hc_{h'}\in\mathfrak m\).
The two families of off-diagonal coefficients are not identically
zero, by irreducibility on \(H_F\). Their valuations have finite
integer minima \(B,C\) with \(B+C\ge1\).
Conjugating by \(\operatorname{diag}(\varpi^{-B},1)\)
retains the diagonal mark and makes every \(b_h\) integral,
some a unit, and every \(c_h\) an element of \(\mathfrak m\).
Consequently in this single basis
\[
 \overline{r(h)}=\begin{pmatrix}1&\bar b_h\\0&\bar\xi(h)\end{pmatrix},
 \qquad h\in H_F,\qquad \bar b_h\ne0\text{ for some }h.
\]
Only stability on \(H_F\) is needed below.

\paragraph{Symplectic induction and the mark.}
Set \(\mu=\epsilon^k\xi\phi\) on \(G_{\mathbb Q}\); it extends
\(\det r\). On \(G_F\), \(\bar\mu=\bar\nu\).
Choose square roots \(a,b\) of the two marked eigenvalues in the
order above, as follows. The first has the unique root reducing
to \(\bar\chi_3(s)=\pm1\) by Hensel's lemma; call it \(a\).
Take \(b=\mu(s)/a\). The determinant identity gives the other
square root, reducing to \(\bar\chi_0(s)\).
Both are units and \(ab=\mu(s)\).

In the ordinary two-coset realization \(I=\operatorname{Ind}_H^{G_{\mathbb Q}}r\)
we may use the matrices
\[
 I(h)=\operatorname{diag}(r(h),r(s^{-1}hs)),\qquad
 I(s)=\begin{pmatrix}0&\operatorname{diag}(a^2,b^2)\\1_2&0\end{pmatrix}.
\]
The coset relation \(I(s)^2=I(s^2)\) and conjugation by \(s\)
verify the representation law. With \(J_2=\left(\begin{smallmatrix}0&1\\-1&0\end{smallmatrix}\right)\),
these matrices preserve \(J_2\oplus abJ_2\) up to multiplier \(\mu\).
For the fixed change of basis
\[
 P=\begin{pmatrix}0&a&0&a\\ b&0&-b&0\\
                   0&-1&0&1\\1&0&1&0\end{pmatrix},\qquad
 J=\begin{pmatrix}0&0&0&1\\0&0&1&0\\0&-1&0&0\\-1&0&0&0\end{pmatrix},
\]
one has \(P^t(J_2\oplus abJ_2)P=-2abJ\).
The determinant of \(P\) is a unit, since \(2,a,b\) are units.
Thus \(R=P^{-1}IP\) has symplectic multiplier \(\mu\) and
\(R(s)=\operatorname{diag}(b,-a,-b,a)\), with precisely the
required residual order.

For \(h\in H\), set \(A=r(h), B=r(s^{-1}hs)\), \(\gamma=b/a\).
Then \(R(h)=M_\gamma(A,B)\), where, with zero-based indices,
\[
2M_\gamma(A,B)=
\begin{pmatrix}
 A_{11}+B_{11}&A_{10}/\gamma-B_{10}&B_{11}-A_{11}&A_{10}/\gamma+B_{10}\\
 \gamma A_{01}-B_{01}&A_{00}+B_{00}&-\gamma A_{01}-B_{01}&A_{00}-B_{00}\\
 B_{11}-A_{11}&-A_{10}/\gamma-B_{10}&A_{11}+B_{11}&-A_{10}/\gamma+B_{10}\\
 \gamma A_{01}+B_{01}&A_{00}-B_{00}&-\gamma A_{01}+B_{01}&A_{00}+B_{00}
\end{pmatrix}.
\]
In particular for \(h\in H_F\) it is integral and reduces to
the prescribed diagonal plus entries in positions
\((1,0),(1,2),(3,0),(3,2)\). This is precisely \(v_V\).
The same holds on the other coset of \(H_F\) in \(G_F\),
by multiplication by the marked diagonal matrix.
At a point with \(\bar A_{01}\ne0\), either
\(\gamma\bar A_{01}-\bar B_{01}\ne0\), or
\(\gamma\bar A_{01}+\bar B_{01}\ne0\); in the latter case both
\((1,2)\) and \((3,0)\) are nonzero. This proves the required
entry alternative.

Since \(\nu/\mu\) is a principal-unit character on \(G_F\),
it has a unique square root \(\delta:G_F\to1+\mathfrak m\).
Indeed squaring is a continuous automorphism of \(1+\mathfrak m\)
for odd \(\ell\), by Hensel's lemma and uniqueness, so the
square root is itself a continuous character. Set
\(Q_r=\delta R|_{G_F}\). It has the exact multiplier \(\nu\)
and unchanged residual support and mark. More explicitly, for
\(d=e+1-k\),
\[
 \delta^2=\epsilon^d\omega^{-1}\psi\phi^{-1}.
 \tag{1}\label{eq:seed-II-twist}
\]
If \(g\in T.\mathrm{fr}(\varnothing).N\), it is inertia away
from \(\ell\) with trivial original residual representation.
Then \(\phi(g)=1\), \(\epsilon(g)=\omega(g)=1\),
\(\psi(g)=1\) by \(\texttt{Ok.kb}\), and \(\xi(g)=1\)
since its reduction is one and it has order prime to \(\ell\).
Both inducing matrices are the identity by the exact tame
inertia assertion, and (1) gives \(\delta(g)=1\).
This verifies the original killing condition, not just its reduction.

\paragraph{Integral local flags.}
For each original \(v\) above \(\ell\), its decomposition group
is in \(H_F\). One inducing factor has a stable line with
residual character \(\bar\xi\), the other with residual
character \(1\). In the marked lattice the first line admits
a primitive vector \((u,1)^t\), \(u\in\mathcal O\), and the
second a primitive vector \((1,w)^t\), \(w\in\mathfrak m\).
Indeed intersect each line with its lattice and take a primitive
vector. On an element of the local group where \(\bar\xi\ne1\),
the upper triangular reduction has distinct diagonal values.
Its \(\bar\xi\)-eigenline has nonzero second coordinate,
whereas its \(1\)-eigenline is the first coordinate line.
This proves the assertions for the two primitive vectors
without changing the marked basis.

Define the lower unipotent matrix
\[
 N(p,q,r,z,f)=
 \begin{pmatrix}1&0&0&0\\p&1&0&0\\q&z&1&0\\r&f&-p&1\end{pmatrix}.
\]
It is symplectic if \(f+pz=q\). Use the following two charts:
\[
\begin{array}{c|ccccc}
 &p&q&r&z&f\\ \hline
 A\text{ has the }\bar\xi\text{-line}
 &\gamma u&-1&\gamma u&
 -2w/(1+\gamma uw)&-(1-\gamma uw)/(1+\gamma uw)\\
 B\text{ has the }\bar\xi\text{-line}
 &-u&1&u&-2w/(\gamma+uw)&(\gamma-uw)/(\gamma+uw).
\end{array}
\]
Their denominators are units. All entries are integral,
\(z\in\mathfrak m\) and \(q=\pm1\), and the symplectic identity
holds. Denote the chosen matrix by \(N_v\).
Here is a verification of the flag calculation. A matrix with
eigenline \((u,1)^t\), line character \(x\), quotient character
\(x'\), has shape
\[
 \begin{pmatrix}x'+u j&(x-x'-u j)u\\j&x-ju\end{pmatrix}.
\]
For \((1,w)^t\) and characters \(y,y'\) its shape is
\[
 \begin{pmatrix}y-iw&i\\(y-y'-iw)w&y'+iw\end{pmatrix}.
\]
Substitution in the displayed formula for \(M_\gamma\)
gives \(N_v^{-1}M_\gamma N_v\) upper triangular with diagonal
\((x,y,y',x')\) in both charts (where \(x\) always denotes
the character of the \(\bar\xi\)-line).
These identities are over the indicated localization of the
integer polynomial ring. In particular no lifting of a residual
flag at the expense of integrality is used. Multiplication by
\(\delta\) retains the flag. Its reductions are ordered as
\((\bar\xi,1,\bar\xi,1)\), and its inertia exponents before the
twist are \(m=(k+t,k,0,-t)\).
The reductions of the lower charts have support in \(q_Q\)
and entry \((2,0)=\pm1\). This proves \(\texttt{InV}\), including
its nonvanishing condition at every original place.

To check the original \(\mathrm{Mu}\) condition, an element
\(g\in\mathrm{Mu}_v\) has an \(\ell\)-power in the closed
commutator subgroup of \(D_v\). The unramified quotient is
\(\widehat{\mathbb Z}\), which has no nontrivial finite-order
element, so \(g\in I_v\). Also \(\epsilon(g)=1\), because
\(\mathbb Z_\ell^\times\) has no nontrivial \(\ell\)-power torsion
for odd \(\ell\). By \(\texttt{Ok.ineL}\), \(\bar\psi(g)=1\).
The depth condition makes \(\psi|_{G_F}\) of order prime to
\(\ell\), so \(\psi(g)=1\), and (1) gives \(\delta(g)=1\).
All four diagonal flag characters are therefore one on
\(\mathrm{Mu}_v\).

Together these give the literal \(\texttt{IsPair}\) and
\(\texttt{Res}\) conditions of the original frame, since
\(v_V\subset n_P\), with continuity over the enlarged DVR.
The local flags are retained as part of this same pair.

\paragraph{Irreducibility on the cyclotomic kernel.}
Let \(U=G_L\cap\ker\epsilon\) (the source's \(T.G'\)).
First \(\phi|_U\ne1\). Otherwise \(\phi|_{G_L}\) factors through
the compact subgroup \(\epsilon(G_L)\subset\mathbb Z_\ell^\times\).
A continuous quadratic character of its pro-\(\ell\) subgroup
is trivial. Hence \(\bar\phi|_{G_L}\) is trivial on
\(\ker\omega\cap G_L\). It permutes the four residual characters
there, since \(\bar\phi=\bar\Psi=\bar\Psi^{-1}\). The precise
\(\texttt{Ok.g3'}\) condition would imply \(\phi|_{G_L}=1\),
contrary to \(\texttt{Ok.g3}\) for \(w_{\rm Ps}\).
Thus \(k_0\not\subset L\), and \(U\) meets the outside coset.

Put \(V=H\cap G_L\), an open subgroup of \(H\).
For each of \(r,r^s\), the identity component of the Zariski
closure of its image on \(V\) acts irreducibly: its inverse image
is open, and the mixed seed is irreducible on every open
subgroup. A connected algebraic subgroup of
\(\mathrm{GL}_2\) acting irreducibly contains \(\mathrm{SL}_2\).
For the unipotent radical fixes a nonzero subspace, which by
normality is invariant, hence the radical is trivial; a
connected reductive subgroup acting irreducibly in dimension two
cannot be a torus, and its derived group contains
\(\mathrm{SL}_2\). The closure of the image of the commutator
of \(V\) therefore contains \(\mathrm{SL}_2\). Since that
commutator lies in \(V\cap U\), the two representations on
\(V\cap U\) are irreducible.

They cannot be isomorphic there. If they were, conjugate one
to agree on \(V\cap U\). This is a normal subgroup of \(V\).
Schur's lemma then shows that on \(V\) the representations
differ by a scalar character; their determinants agree,
so this character has square one. Hence their characteristic
polynomials agree on every square from \(V\).
Choose a nonzero power of the original \(\ell\)-inertia element
with cyclotomic value of infinite order (from the original
good datum), lying in \(V\); it exists by finite index.
On its square the two inertia eigenvalue multisets are
\[
 \{\epsilon^{k+t},\epsilon^{-t}\},\qquad
 \{\epsilon^k,1\}.
\]
They differ, as \(\epsilon\) has infinite order and \(k,t>0\).
This contradiction proves nonisomorphism. Now the restriction
of the induction to \(V\cap U\) is the direct sum of these two
nonisomorphic simple modules. A proper nonzero invariant
subspace for \(U\) would have to be one of the summands,
but an element of \(U\) outside \(H\) interchanges them.
Thus \(R|_U\), and hence \(Q_r|_U\), is absolutely irreducible.

\paragraph{The cuspidal attachment and its local bounds.}
Write \(M=k_0L\). The extension \(M/k_0\) is abelian, and
\(M/L\) is quadratic. Take the Hilbert cuspidal representation
in the source normalization of Section~\ref{sec:mixed-hilbert-seed}
and base change it to \(M\) along a cyclic tower. It remains
cuspidal: at each step a nontrivial self-twist (the obstruction
in cyclic base change) would, by the unramified characteristic
polynomials, Chebotarev and semisimplicity, give a nontrivial
self-twist of the corresponding Galois representation. On the
kernel of the twisting character that representation remains
irreducible, so Schur's lemma forces the character to be trivial.
The base change to \(M\) is not invariant under
\(\operatorname{Gal}(M/L)\); the two embeddings above each
embedding of \(L\) restrict to different embeddings of \(k_0\)
and have distinct weight multisets.

Use cyclic automorphic induction of this cuspidal representation
to \(L\). It is cuspidal by the noninvariance just proved.
These applications use Arthur--Clozel's cyclic base change and
automorphic induction, or Henniart, \emph{Induction automorphe
globale pour les corps de nombres}, Bull.\ Soc.\ Math.\ France
140 (2012), Theorems 1--3. The finite local compatibility is
Theorems 4--5 there, and the infinite-place compatibility is
the concluding remark of Section 3.7. All these results may
be applied after unitary normalization and then untwisted.

Use geometric Frobenius for the attachment: the normalized
rank-two parameters are the eigenvalues of geometric Frobenius
divided by \(q^{1/2}\). The normalized rank-four parameters
of the induction therefore require the common extra twist
\(|\det|\), dividing its parameters by \(q\).
At a split unramified prime the induced characteristic polynomial
is the product of the two rank-two polynomials; at an inert
unramified prime it is \(\det(X^2-r(g^2))\). The second formula
follows by the two-coset realization and accounts for both square
roots. These identities verify the rank-four attachment, with
the factor \(q^{3/2}\) in the source's definition.
At infinity the multiset in this same normalization is
\(\{k+t,k,0,-t\}+3/2\), which is regular algebraic.
Denote the resulting cuspidal representation by \(\Pi\).

We also need the full finite-place window, not merely its
unramified version. The starting Hilbert packet has holomorphic
weights \(k+2t+1,k+1\), at least two and of the same parity.
By Blasius, \emph{Hilbert modular forms and the Ramanujan
conjecture} (arXiv:math/0511007), Theorem 1, its unitary
normalization is tempered at every finite place.
Its local types are principal series with distinct inertia
characters at its conductor primes, and spherical elsewhere,
including at \(\ell\), by the mixed Hilbert construction.
The extension \(M/L\) is unramified at finite places by
\(\texttt{Two.unr}\) and split above \(2,\ell\).
Cyclic local base change preserves the principal-series
parameter description. At split places the local induction
takes the sum of two such parameters; at nonsplit places the
starting parameter is unramified (the conductor is over \(2\))
and induction along an unramified quadratic extension is again
a sum of unramified characters. All are unitary after the
single common global normalization. Thus after the final norm
twist there is one real number \(z_0\) such that every local
degree-one support character at a finite place of residue size
\(q\) has absolute value \(q^{z_0}\) on a uniformizer.
This lies strictly between \(q^{z_0-1/2}\) and
\(q^{z_0+1/2}\). This is the literal \(\texttt{AllWin}\)
condition. It holds for the ramified places as well, and uses
no unknown rank-four Ramanujan assertion.

Finally take \(\kappa=\delta^{-1}|_{G_L}\).
Then \((Q_r|_{G_L})\otimes\kappa=R|_{G_L}\) is attached
to \(\Pi\). By depth, \(\psi|_{G_F}\) has finite order dividing
some positive integer \(m_0\) prime to \(\ell\). Set
\[
 h_i=2m_i+d,\qquad a_0=-d,\qquad N_0=2(\ell-1)m_0 .
\]
On the open subgroup of \(G_F\) killing \(\omega,\psi\),
equation (1) and the exact inertia characters prove
\(\texttt{Wts}(h)\), since \(\phi\) is trivial on local inertia.
The consecutive gaps of \(h\) are twice integers at least two;
\(a_0+h_0=2(k+t)\) is even. Moreover
\((\kappa^2\epsilon^{-a_0})^{N_0}=1\) on \(G_L\),
again by (1). The twist is continuous and \(N_0>0\).
These are all the clauses of \(\texttt{IsSeedAlt}\).
This completes the construction of the original Case II
seed pair over one finite coefficient extension.
\end{proof}

The construction separates three issues: a cuspidal seed in rank
two, the marked integral symplectic induction, and the transport
of automorphy from the resulting rank-four seed to arbitrary
points. The theorem proves the second and the seed's own
automorphic attachment. The last issue still requires the
residual linkage and the rank-four potential-modularity ideal
with its arithmetic point and characteristic-prime propagation
properties. No dimension estimate on that ideal is a consequence
of the existence of the seed.
